\chapter{Diskussion}\label{ch:Diskussion}

\section{Interpretation der Ergebnisse}\label{ch:Diskussion:Interpretation der Ergebnisse}

\section{Einordnung gegenüber publizierten Ergebnissen}\label{ch:Diskussion:Einordnung gegenüber publizierten Ergebnissen}

\section{Threats to Validity}\label{ch:Diskussion:Threats to Validity}

\section{Belastbarkeit der Ergebnisse}\label{ch:Diskussion:Belastbarkeit der Ergebnisse}

\section{Implikationen und Handlungsempfehlungen}\label{ch:Diskussion:Implikationen und Handlungsempfehlungen}

\section{Offene Fragen}\label{ch:Diskussion:Offene Fragen}

\controlquestions{
Dieses Kapitel in einer Abschlussarbeit ist zentral, um die Ergebnisse einzuordnen, kritisch zu reflektieren und sie in einen größeren wissenschaftlichen Kontext zu stellen. 
Hier wird gezeigt, dass die forschende Person Arbeitsergebnisse nicht nur präsentieren, sondern auch durchdenken,  qualifiziert einordnen und bewerten kann.
\begin{itemize}
    \item Interpretation und Bewertung der Ergebnisse
        \begin{enumerate}
            \item Werden die Ergebnisse kritisch interpretiert – auch in Hinblick auf die gestellten Anforderungen oder Hypothesen?
            \item Wird erklärt, was die Ergebnisse bedeuten – im Lichte der Zielsetzung, Hypothesen oder Erwartungen?
            \item Wird zwischen objektiven Befunden und subjektiver Interpretation (persönliche Bewertung) klar unterschieden?
            \item Wird erkennbar, ob und wie gut das entwickelte System/Konzept den Anforderungen gerecht wird?
            \item Werden die zentralen Ergebnisse sinnvoll zusammengefasst – ohne sie einfach zu wiederholen?
            \item Wurde klar zwischen Korrelation und Kausalität unterschieden?
        \end{enumerate}
    \item Reflexion von Unsicherheiten und Einschränkungen
        \begin{enumerate}[resume]
            \item Werden mögliche Fehlerquellen, Unsicherheiten oder Einschränkungen der Ergebnisse reflektiert?
            \item Wurden unerwartete oder widersprüchliche Ergebnisse angesprochen und erklärt?
            \item Wird die Zuverlässigkeit und Aussagekraft der Ergebnisse eingeschätzt?
            \item Sind methodische Schwächen oder Einschränkungen der Arbeit ehrlich und sachlich reflektiert?
            \item Ist transparent, was methodisch nicht optimal gelöst werden konnte – und warum?
            \item Wird die Einschätzung der Bedrohungen sachlich, reflektiert und ohne Beschönigung dargestellt?
            \item Wird klar, wie stark die Ergebnisse trotz der \emph{Threats of Validity} belastbar sind?
        \end{enumerate}
    \item Einordnung im wissenschaftlichen Kontext
        \begin{enumerate}[resume]
            \item Werden die Ergebnisse in Bezug zu bestehender Literatur oder Theorien (vgl.~Kapitel \ref{ch:Verwandte Arbeiten}) gesetzt?
            \item Wird erläutert, inwiefern die Ergebnisse bestehende Ansätze stützen, erweitern oder widersprechen?
            \item Werden relevante wissenschaftliche Kontroversen oder alternative Erklärungen aufgegriffen?
            \setcounter{lastenumcounter}{\value{enumi}}
        \end{enumerate}
\end{itemize}
}
\controlquestionscontinue{
\begin{itemize}
    \item Praktische und theoretische Implikationen
        \begin{enumerate}\setcounter{enumi}{\value{lastenumcounter}}
            \item Ergeben sich aus den Ergebnissen Handlungsempfehlungen, Optimierungspotenzial oder weitere Fragestellungen?
            \item Welche praktischen, wissenschaftlichen oder gesellschaftlichen Konsequenzen lassen sich aus den Ergebnissen ableiten?
            \item Gibt es Handlungsempfehlungen oder Hinweise für die Anwendung in der Praxis und/oder weiteren Forschung?
        \end{enumerate}
    \item Überleitung zu Ausblick und zukünftige Arbeiten
        \begin{enumerate}[resume]
            \item Welche offenen Fragen bleiben bestehen?
            \item Welche Ansätze zur Weiterentwicklung oder für zukünftige Forschung ergeben sich?
            \item Wird die Diskussion schlüssig abgeschlossen und ein Übergang zum Fazit (vgl.~Kapitel \ref{ch:Zusammenfassung und zukünftige Arbeiten}) geschaffen?
        \end{enumerate}
    \item Threats to Validity
        \begin{enumerate}[resume]
            \item Gab es Störfaktoren, die das Ergebnis beeinflusst haben könnten (z.\,B.~Lernkurveneffekte, unbeachtete Variablen)?
            \item Wurden potenzielle Bias (z.\,B.~durch Auswahl der Probanden, Messinstrumente) reflektiert?
            \item Ist die Stichprobe (z.\,B.~Personen, Daten, Testfälle) repräsentativ genug für eine Verallgemeinerung?
            \item Wird erklärt, in welchen Kontexten die Ergebnisse nicht ohne Weiteres gelten?
            \item Sind die Rahmenbedingungen der Untersuchung möglicherweise so speziell, dass sie die Übertragbarkeit einschränken?
        \end{enumerate}
\end{itemize}
}


